% Abhakseite „Das kann ich“ – Zone nur im Gesamt (\mitzone)
%% TODO Ich-kann-Sätze: Platzhalter wie in den Aufgabendateien
\begin{abhakseite}
\ifdefined\mitzone
\abhakgruppe{Kennst du schon}
\abhak{1}{Flächen- und Umfangsformeln von Rechteck, Dreieck und Trapez, Satz des Pythagoras für Streckenlängen}
\abhak{2}{Punkte (x | f(x)) im Koordinatensystem lesen und Figuren einzeichnen, Funktionswerte berechnen}
\abhak{3}{Terme ausmultiplizieren und zusammenfassen}
\abhak{4}{Ableitungen bilden: Potenz-, Faktor- und Summenregel; GK und LK zusätzlich Produkt- und Kettenregel für e-Funktions-Produkte}
\abhak{5}{Gleichungen lösen: linear, quadratisch mit Lösungsformel, biquadratisch mit Substitution, Ausklammern und Nullprodukt}
\abhak{6}{Extrempunkte über notwendige und hinreichende Bedingung, Randvergleich auf einem Intervall}
\abhak{7}{Flächen- und Umfangsformeln von Rechteck, Dreieck und Trapez, Satz des Pythagoras für Streckenlängen – Fehler finden}
\abhak{8}{Flächen- und Umfangsformeln von Rechteck, Dreieck und Trapez, Satz des Pythagoras für Streckenlängen – selbst rechnen}
\fi
\abhakgruppe{Einheit 1 · Figur und Term}
\abhak{9}{Figur und Term}
\abhak{10}{Figur und Term – Fehler finden, Begründen, Darstellungswechsel}
\abhakgruppe{Einheit 2 · Zielfunktion aus Haupt- und Nebenbedingung}
\abhak{11}{Zielfunktion aufstellen}
\abhak{12}{Zielfunktion aufstellen – Fehler finden, Begründen, Anwendung}
\abhakgruppe{Einheit 3 · Maximum bestimmen und deuten}
\abhak{13}{Maximum bestimmen}
\abhak{14}{Kleinsten Abstand eines Punktes zu einem Graphen über die Abstandsfunktion bestimmen}
\abhak{15}{Maximum bestimmen – Fehler finden, Begründen, Anwendung}
\end{abhakseite}
